\scanpage{6}
\mspage{4}
Je me rends compte que mes notations \unsure{ici ne|ne me} sont pas commodes pour la dualité, je vais les changer, en $u_i$ pour les flèches entre objets $\{i\}$ avec $i\geq0$, \deleted{\unsure{et}} et $v_i$ la flèche symétrique de $u_i$ p.r. à l'origine. (Le groupe des autom. de $\Phi_n$ est réduit au groupe $\mathbb Z_{/2}$ à deux \unsure{élémts|éléments}, le seul automorphisme non trivial étant la symétrie $\{\deleted{\unsure{x}}i\}\longrightarrow\{-i\}$. Donc nos \deleted{\unsure{diagrammes}} chemins seront notés désormais
\diagram{\begin{tikzcd}[ampersand replacement=\&,column sep=.9em]
(*)\quad A_{-2}\arrow[r,"v_1"] \& A_{-1} \& A_0\arrow[l,"v_0"']\arrow[r,dash,"u_0"] \& A_1 \& A_2\arrow[l,"u_1"']\arrow[r] \& \cdots
\end{tikzcd}}
(Ce n'est pas non plus \unsure{enthousiasmant}\ldots)

\deleted{On a deux \unsure{mors} de la} Soit \deleted{$\unsure{\varepsilon|\mathcal E}$}
\[
\varepsilon=\{-1,+1\}
\]
l'ens ordonné \uline{discret}, réduit aux deux \unsure{élémts|éléments} $-1,+1$, et définissons des foncteurs d'inclusion (pl. fidèles)
\[
\begin{aligned}
\alpha_n:\varepsilon\longrightarrow\Phi_n\qquad &\alpha_n(-1)=-n\\
&\alpha_n(1)=n\quad,
\end{aligned}
\]
ces foncteurs sont compatibles avec les $\rho_{n,m}$, d'où
\[
\alpha_\infty:\varepsilon\longrightarrow\underset{\underset{\mathrm{Pro\,Cat}}{\rotatebox{90}{$\in$}}}{\mathbb{\Phi}}\quad,
\]
d'où pour tout $X$ un foncteur
\[
\unsure{\mathcal L|\ell|\bigl(}\qquad\underline{\mathrm{Hom}}(\mathbb{\Phi},X)\longrightarrow\underline{\mathrm{Hom}}(\varepsilon,X)\simeq X\times X
\]
i.e.

% archive sheet number: 5
